\section{Soviet Disarray: Rival Bureaus and the N1 Failures}

\subsection{Korolev, Chelomei, and Divided Management}
The Soviet crewed lunar effort suffered from deep organizational rivalry between competing design bureaus rather than the unified direction that characterized Apollo. Sergei Korolev's OKB-1 and Vladimir Chelomei's bureau pursued separate approaches to reaching the Moon, each backed by different patrons within the Soviet leadership. This competition divided scarce resources, engineering talent, and political attention across parallel programs. Instead of concentrating national effort behind a single architecture, the Soviet system dispersed it, undermining the coherence that a project of such complexity demanded.

Korolev, the architect of Sputnik and Gagarin's flight, favored his own launch vehicle and mission design, while Chelomei advanced rival concepts supported by his own influential connections. The lack of a single empowered authority to adjudicate between them meant that decisions were repeatedly contested and delayed. Funding shifted with political favor rather than technical merit, and overlapping work wasted effort that a centralized program would have directed efficiently. The bureaucratic struggle consumed energy that should have been spent solving the formidable engineering problems of lunar flight.

The death of Sergei Korolev in January 1966 removed the program's most capable and determined leader at a critical moment. Korolev had possessed the prestige and willpower to push his projects forward despite obstacles, and his absence left a vacuum that no successor filled with equal authority. Vasily Mishin, who took over OKB-1, lacked Korolev's political influence and command of the sprawling effort. The loss of this driving figure contributed significantly to the delays and disorganization that followed in the lunar program.

Soviet secrecy compounded the structural problems by limiting coordination and obscuring accountability. The program operated under tight security that prevented open debate and made failures difficult to confront and correct. Fear of public embarrassment encouraged officials to conceal setbacks rather than address them transparently, the opposite of the visible commitment that drove Apollo. This culture of concealment, combined with divided management, deprived the Soviet effort of the corrective feedback that a complex technological program requires to succeed.

The contrast with the American approach was stark and consequential. Where NASA concentrated authority and resources behind a single publicly announced goal, the Soviet Union fractured its effort among rival bureaus competing for influence. The absence of unified management meant that strategic decisions were slow, contested, and frequently reversed. Historians continue to debate whether a consolidated Soviet program under a single leader could have reached the Moon first, but the divided structure that actually existed proved a decisive handicap in the race. \cite{harford1997} \cite{siddiqi2010}

\subsection{The Four Failed N1 Launches}
At the heart of the Soviet crewed lunar plan stood the N1, an enormous heavy-lift rocket intended to carry cosmonauts toward the Moon. The vehicle was designed to rival the American Saturn V in lifting capacity, but it embodied a fundamentally different and riskier engineering philosophy. Its first stage clustered thirty NK-15 engines whose combined thrust was meant to overcome the rocket's great mass. This unprecedented arrangement created daunting challenges in plumbing, control, and coordination that the program never fully mastered.

The decision to use thirty engines reflected the limited thrust of available Soviet engines compared with the massive American powerplants. While the Saturn V relied on five large, well-tested engines in its first stage, the N1 distributed its load across many smaller units. Managing the simultaneous operation of so many engines, with their interconnected fuel lines and control systems, proved extraordinarily difficult. Any failure among the cluster could cascade into catastrophe, and the complexity overwhelmed the testing and quality control resources available to the program.

All four test launches of the N1 between 1969 and 1972 ended in failure, none reaching orbit. The first attempt in February 1969 failed shortly after liftoff, and the second in July 1969 fell back onto its launch complex and exploded with devastating force. That second failure destroyed the pad and set the program back severely, occurring just before Apollo 11 reached the Moon. Subsequent launches in 1971 and 1972 likewise failed, confirming that the fundamental problems remained unsolved despite repeated effort.

A crucial weakness lay in the inability to test the assembled first stage on the ground before flight. Unlike the American practice of extensive static firing, the Soviet program flew complete stages without comprehensive prior testing, relying on the launches themselves to reveal problems. This approach meant that each failure destroyed an entire vehicle and the data it might have provided over many flights. The combination of inadequate testing and immense complexity made the path to a reliable N1 painfully slow and expensive.

The repeated N1 failures effectively doomed the Soviet crewed lunar effort. By the time engineers had begun to understand the engine and control problems, the United States had already landed astronauts on the Moon, eliminating the prestige that justified the costly program. Soviet leaders quietly canceled the N1 in the mid-1970s, concealing its existence for decades afterward. The rocket's failure stands as the technical culmination of the organizational weaknesses that had hampered the Soviet program from its divided beginnings. \cite{harford1997} \cite{siddiqi2010}

\begin{figure}[htbp]
\centering
\includegraphics[width=0.88\linewidth]{figures/ch04_lunar_surface.jpg}
\caption{Apollo 11 Lunar Module Eagle on the Moon surface.}
\label{fig:ch_apollo_11_lunar_module_eagle_on_}
\end{figure}

\begin{figure}[htbp]
\centering
\includegraphics[width=0.92\linewidth]{figures/mission_timeline.pdf}
\caption{Moon Race milestones chart generated with Python (matplotlib).}
\label{fig:timeline}
\end{figure}

\bigskip
\begin{otherlanguage}{hebrew}
\subsection*{סיכום הפרק}

התחרות בין לשכות התכנון של קורוליוב וצ'לומיי פיצלה את המשאבים ופגעה בתוכנית הירח הסובייטית המאוישת. רקטת ה-{\beginL N1\endL} הכבדה נכשלה בכל ארבעת שיגורי הניסוי שלה בין {\beginL 1969\endL} ל-{\beginL 1972\endL}.

\end{otherlanguage}

\clearpage
